% §1 Introduction --- written last (Q-0292, 2026-10-04), from the finished sections 2-5 and the
% appendices, on Paper V's model (spatial-hemigroup-scale-space, paper-b/01-introduction.tex): the
% reader's starting point (Gaussian scale space and shape-adapted smoothing), the question, the
% setting of the line papers, the answer in outline with the two tables, what is classical and what
% is new, the contributions by statement number, the scope, the plan, and 1.1. The draft's §1
% (draft/spatial-hemigroup-affine.md) was the starting text. Every sentence that restates a result
% was compared with its statement (WRITING.md § 5, "a summary drops hypotheses"). Citations: only
% those of sections 2-5, at the anchors used there (lindeberg1997shape by section only, the stored
% copy being the report version). Nothing here is shared with the blueprint.
%
% 1.1 owns the trust base (article-kit PROCESS.md § 7, WRITING.md § 4): the state of the
% formalization (the Lean statement skeleton is in progress, Q-0281; no \leanok), the cited facts
% with the statements that rest on them (the ledger's **Blueprint:** lists, blueprint/AXIOMS.md, and
% the cited-fact sentences of sections 3-4), the sketched proofs, and the computed remarks and
% figures.
%
% After the blind review (Q-0293; records/affine/REVIEW-paper-sections-1-6.md): the summaries
% narrowed to their statements (the positive definite base, 0 < alpha, gamma < 2, "other axioms",
% three exceptions to "proved in the text"); the cited-fact map exhaustive by section from the
% ledger's Blueprint lines; prior art kept to the sources read (meerschaert2001limit added, its
% reading note in records/affine/), with TODO(source) comments where a source not held belongs.
%
% Q-0294: schoenberg1938metric, bryc1995normal and holmes1982operator read at the page and cited
% (schoenberg1938positive read, not cited; READING notes in records/affine/). The TODO(source)
% comments left are Hudson-Mason (1981) and Jurek (1992; cited as 1991 in Jurek–Mason), none held;
% the Gneiting TODO went with the Schoenberg comparison to the isotropic paper.
%
% Q-0307 (the split at §3 | §4, adr/0001): the introduction of the affine paper alone. It leads
% with shape adaptation, as the whole-paper review asked (records/affine/REVIEW-paper-whole.md, flag
% 50): the Lindeberg-Garding practice and the answer for it (prop:shape-adapted-apertures, the
% single article's fifth "answer" paragraph, moved up) and the rational symbols
% (prop:rational-symbols, from the single article's section 6, "What is elsewhere") come second.
% The isotropic theory is the isotropic paper's: its "answer" paragraph and contribution 2 went
% there, and the corner paragraph here keeps one sentence of it. The rest is the single article's
% text, moved; §1.1 keeps the cited facts this paper spends (A3, A5, A6) and the sections that spend
% them. The fixes of the review (flags 1-7, 27-29, 42, 48-49) are the next item's, not this one's.
%
% Q-0309 (the review's fixes; records/affine/REVIEW-paper-whole-fixes-affine.md): one home per
% comparison and disclaimer (flags 27-32, 42). "What is classical and what is new" is one positioning
% paragraph with pointers: Lindeberg's hypotheses live in rem:lindeberg, Kelker and Herschel-Maxwell in
% rem:kelker, the morphological and alpha-scale-space axiomatics in 3.5, the cone-field frame in 4.4
% ("Prior art"), Jurek-Mason, Meerschaert-Scheffler and Holmes-Hudson-Mason in 4.6 ("Prior art"),
% Whittle and Lindgren-Rue-Lindstrom after prop:affine-orbits. "Scope" is the one home of the
% kernels-only disclaimer (with the operator reading of the rigidity, from section 5's former
% paragraph "Stated on kernels" and Appendix A); singularity theory is section 5's "Not treated".
% Table 1 (the corners) is dropped: the corners are cells of tab:landscape (flag 34).
% Novelty hedges (fix group 5, records/affine/PRIOR-ART-novelty-claims.md): the contributions list
% carries "against the sources read" for K_alpha, the budget and the step cone (flag 49) and for the
% single-step form (Kelker, Fang-Kotz-Ng, Bryc, Cambanis-Huang-Simons; flag 46); the stable orders are
% claimed only as the orders of the stable orbits (flag 47). The TODO(source) on Hudson-Mason and
% Jurek 1992 is closed: both are read and cited in 4.6.
%
% Q-0314: Pedersen-Sato 2003 read and cited (the canonical PSD(d)-parameter semigroup); "Scope" claims
% the operator form (lem:kernels-to-operators, thm:affine-rigidity-operators; no cited fact added, the
% converse needing L^1 convergence only, which Gaussian smoothing gives); the ray and Loewner orders
% taught where first named; three summaries narrowed to their statements
% (records/affine/SCOPE-AUDIT-affine.md).
%
% Q-0377 (the formalization closed: 125 of 125 blueprint nodes \leanok, main at dfafe99): 1.1 is
% rewritten as "What is machine-checked" on Paper V's model, with the headline #print axioms block
% and the trust-base markers `linkage check` reads; the contributions end with one clause on the
% formalization. The fidelity review (Q-0376) is named as pending in 1.1.

% The writing-rules pass (group s1, 2026-10-14; records/affine/edits_wr_s1.py): section 1
% restructured on the released isotropic paper's model under the author's rules 2, 6 and 12: the
% reader's starting point, the axioms, the change to the vocabulary of probability, the (G, W)
% pair, the questions, the answers at the abstract's level with the statement of each, scope, the
% two relation paragraphs, the plan. The contributions list is folded into the answers and the
% plan; "corner", "row", "column", "landscape" and "frame" (for a setting) are gone from the
% section; the claims of novelty against the sources read are gone (the author's direction of
% 2026-10-12). Every sentence that restates a result was compared with its statement.
\section{Introduction}\label{sec:introduction}

A scale-space representation measures an image through an aperture whose size is the scale. A
measurement is convolution with a probability kernel, and the record at a coarser scale is computed
from the record at a finer one by a further convolution, the \emph{cascade}. In dimension $d \ge 2$
the aperture also has a shape. The affine Gaussian scale space of Lindeberg and G\aa{}rding makes the
shape a parameter. Its kernel at the matrix $\Sigma$ is the Gaussian with covariance $\Sigma$, the
image is smoothed once for each $\Sigma$, and an invertible linear change $B$ of the image coordinates
moves the parameter by $\Sigma \mapsto B\Sigma B^{\tp}$ (\sscite{lindeberg1994shape},
eqs.~(12)--(16), pp.~393--394; \sscite{lindeberg1997shape}, \S4). Its use is shape adaptation. The iteration of Lindeberg and
G\aa{}rding fits the shape of the aperture at a point to the second-moment matrix of the image there,
so that what is measured is the same for every affine image of a pattern (\sscite{lindeberg1994shape},
pp.~395--396; \sscite{lindeberg1997shape}, \S\S5.2 and~6.2). The nonlinear filter of Nitzberg and Shiota likewise
shapes a Gaussian kernel at each point by the gradients of the image around it, to enhance edges and corners
(\sscite{nitzberg1992nonlinear}, eqs.~(3), (3') and~(4), p.~827). The Gaussian is also what the classical axioms of scale
space select (``Relation to scale-space theory'' below). This paper asks which affine scale spaces
exist under other axioms, incomparable with the classical ones (Remark~\ref{rem:lindeberg}), and what
each of them can do for shape adaptation.

% Pass 6, the author's review of section 1 (2026-10-26; records/affine/REVIEW-author-paper-affine.md,
% comment 1): the question is posed as the positioning brief posed it (adr/0006, items 2-4; the
% brief itself is deleted). First why the question arises (the Gaussian cascades along every widening, so
% neither the kernel nor the order is a choice; under the hemigroup axioms both are), then the two
% choices defined with what each one is (Sigma carries size and shape, the scale being its part
% along the dilations; G with the translations is the affine group the family respects; W is the
% set of stages), then the main result, the classification of Table tab:landscape, read in three
% parts: the general linear group with the Loewner order (with shape adaptation), the general
% linear group from the rays to the Loewner order (the headline: the price of a non-Gaussian kernel
% kernel by kernel), and the smaller groups (the rigidity is a property of the group). The four
% questions and the four "answers" are gone; the text of the answers is kept, under the three
% headings. Every new sentence that restates a result was compared with its statement:
% prop:ray-families, thm:affine-rigidity, prop:stable-step-cone(1)-(2) (the stable order is the
% step set of the stable orbit), prop:matern-steps, lem:gw-kernels-from-origin(1),
% prop:similarity-stable-family(3), prop:fixed-frame-families with section 4.9, "Frames".
% Round 2 of the external review (the presentation review, Gemini; the author's decision D6,
% 2026-11-03, records/affine/PLAN-review-response-paper.md): "What is an affine scale space?" and
% "The two choices" come before the axioms, the earlier papers and the vocabulary of probability, so
% that the organizing idea comes first. The first paragraph uses no term that the vocabulary
% paragraph teaches: the orbit's hypotheses are "of the kind the axioms below allow", and the
% stable and Matern sentences, with the name "affine orbit", move to the end of the vocabulary
% paragraph. "The axioms" is a new paragraph heading.
\paragraph{What is an affine scale space?} In the affine Gaussian scale space the cascade needs no
choice. The kernel is the Gaussian, and the Gaussian of covariance $\Sigma'$ is the
Gaussian of covariance $\Sigma$ convolved with the Gaussian of covariance $\Sigma' - \Sigma$ whenever
that difference is positive semidefinite. So every widening of the aperture is a further smoothing,
and there is nothing to choose. Under other axioms the kernel need not be the Gaussian, and then the
cascade becomes a choice. A rotation-invariant kernel of the kind the axioms below allow, a
self-decomposable one in the sense defined below, can be carried to every elliptical aperture by the
linear maps of the image. For $d \ge 2$ and a kernel other
than the Gaussian, the family it gives cascades along every zoom, a change of size that keeps the
shape, and never along every change of shape (Propositions~\ref{prop:ray-families}
and~\ref{prop:affine-orbits}(3)). So the name ``affine scale space'' leaves two things open: which changes of aperture must be further smoothings, and which
linear maps of the image the family must respect. Several answers are reasonable. The changes of
aperture can be the zooms alone or every widening, and the maps all invertible linear maps, the
rotations with the dilations, or the independent rescalings of the axes of a sampling grid. The
name already has more than one sense in the literature. Lindeberg's affine scale space is the family
of Gaussians of every covariance matrix, closed under the invertible linear maps of the plane
(\sscite{lindeberg1994scale}, (15.17), p.~388). Weickert's affine Gaussian scale space is the
Gaussians of covariance $tD$ for one fixed matrix $D$ (\sscite{weickert1998anisotropic}, \S1.2.6,
p.~13). The affine invariant scale spaces of Alvarez, Guichard, Lions and Morel and of Sapiro and
Tannenbaum are nonlinear evolutions with a scalar time (\sscite{alvarez1993axioms}, Thm.~5, p.~229;
\sscite{sapiro1993affine}, \S6, p.~32). This paper takes the name in Lindeberg's sense with the
kernel left free, and its question is a classification over the two choices.

\paragraph{The two choices.} In dimension $d$ an aperture is indexed by a positive semidefinite matrix
$\Sigma$, proportional to its covariance in the Gaussian case, which fixes its size and its shape
together. A scalar scale is only one part of this parameter. A dilation $x \mapsto \lambda x$
of the image, $\lambda > 0$, moves $\Sigma$ to $\lambda^2\Sigma$, along its \emph{ray}, and changes
the size of the aperture alone. The dilations form a subgroup of the invertible linear maps, and a
general map $A$ moves $\Sigma$ to $A\Sigma A^{\tp}$, changing the shape as well. A family has a kernel
$\mu_{\Sigma,\Sigma'}$ for each \emph{stage}, from an aperture $\Sigma$ to a wider one $\Sigma'$. The
two choices fix which stages there are and how they transform.
\begin{itemize}
\item The \emph{covariance group} $G$ is a closed group of invertible linear maps of $\RR^d$ that
contains the dilations. The family is covariant under it: for $A \in G$ the kernel of the stage from
$A\Sigma A^{\tp}$ to $A\Sigma'A^{\tp}$ is the image under $A$ of the kernel of the stage from $\Sigma$
to $\Sigma'$. Every family commutes with the translations, so the affine maps of the image that it
respects are those whose linear part lies in $G$. Covariance lets a measurement made at one shape be
carried to another by the image's own change of coordinates.
\item The \emph{cascade order} $W$ is a partial order $\le_W$ on the positive semidefinite matrices
that is invariant under $G$ and relates each $\Sigma$ to its multiples $c\,\Sigma$, $c \ge 1$, and $0$
to every matrix. Its pairs $\Sigma \le_W \Sigma'$ are the stages
of the family, and along them the stages compose,
$\mu_{\Sigma,\Sigma''} = \mu_{\Sigma',\Sigma''} * \mu_{\Sigma,\Sigma'}$ for
$\Sigma \le_W \Sigma' \le_W \Sigma''$. So $W$ is the set of the changes of aperture along which the
record at the coarser aperture is a further smoothing of the record at the finer one. Off $W$ the family
has no stage, so a change of aperture need not be a further smoothing. A change with no increment that
is a probability measure is a \emph{re-measurement} of the original image.
\end{itemize}
Two orders bound the ones studied here. The \emph{ray order} relates $\Sigma$ only to its multiples
$c\,\Sigma$, $c \ge 1$, and $0$ to every matrix, so a stage grows the aperture and keeps its shape.
The \emph{Loewner order} relates $\Sigma$ to $\Sigma'$ whenever $\Sigma' - \Sigma$ is positive
semidefinite, so a stage shrinks the aperture in no direction. One definition covers every pair, the
$(G, W)$ cascade scheme of Definition~\ref{def:gw-cascade-family}: kernels $\mu_{\Sigma,\Sigma'}$,
indexed by the pairs $\Sigma \le_W \Sigma'$, that cascade along $W$, depend continuously on the pair,
are not trivial from $0$ to $I$, and are covariant under $G$. The aperture at $\Sigma$ is the kernel
$\mu_\Sigma := \mu_{0,\Sigma}$ of the stage from $0$, the kernel \emph{from scale zero}. The affine Gaussian scale space is a family of the
largest group, the general linear group $\GL(d)$, along the largest of these orders, the Loewner
order.

\paragraph{The axioms.} Scale-space axioms usually require the cascade to be a semigroup in one scale parameter, as those of
Pauwels, Van Gool, Fiddelaers and Moons do \sscite{pauwels1995extended}, so that one parameter indexes
the stages and the apertures alike. The axioms used here require only that the stages compose. With one
scale parameter, as on the line, and $\mu_{s,t}$ the kernel of the stage from scale $s$ to scale $t$,
this is
$\mu_{r,t} = \mu_{s,t} * \mu_{r,s}$ for $r \le s \le t$, and how a stage depends on its two scales is
left free. Such a family is a \emph{convolution hemigroup} (\sscite{hazod2001stable}, Def.~2.14.16,
p.~380). On the line the axioms are those of the author's paper on the line
\sscite{fagerstrom2026spatial}, the \emph{line paper} below (Definition~\ref{def:line-family}). They
require positive, mass-preserving stages that cascade, commute with translations and the reflection
and depend continuously on the scales, and a dilation of the signal may only relabel the scales. Its
sequel \sscite{fagerstrom2026cone}, the \emph{cone paper}, describes the line's kernels further,
and the \emph{isotropic paper} \sscite{fagerstrom2026isotropic} carries the theory to $\RR^d$ for
apertures that have a size and no shape. What is used from the three is restated in
\S\ref{sec:line}.

A kernel is a probability measure, and the classification on the line, like the study of the
apertures that follows, rests on the theory of such measures in probability. So from here on the
paper uses the vocabulary of that theory: a kernel is a \emph{law}, and its Fourier transform, the
transfer function of the smoothing, is the characteristic function of the law, its \emph{transform}
for short. Scale, stage and cascade keep their meaning, and ``kernel'' stays in use for a law in its
role as an aperture or as a stage. On the line the families are known (Theorem~\ref{thm:line-classification}). In
a suitable parametrization of the scales, the kernels from scale zero are the dilations of one symmetric \emph{self-decomposable} law. For every
$0 < c < 1$ such a law is the convolution of its own image under $x \mapsto cx$ with another law
(\sscite{sato1999levy}, Def.~15.1, p.~90). In a cascade this is the relation between two scales: the
kernel at the finer scale is a contracted copy of the kernel at the coarser one, and the other law is
the stage between them. Every such law other than the point mass occurs. Among them are the symmetric
stable laws, with transforms $e^{-a|\omega|^p}$, $a > 0$, $0 < p \le 2$, the Gaussian at $p = 2$
(\sscite{sato1999levy}, Thm.~14.14, p.~86), and the \Matern{} laws, with transforms
$(1 + \omega^2)^{-\gamma}$, $\gamma > 0$, which carry the name of the \Matern{} covariances of spatial
statistics (\sscite{matern1960spatial}, (2.4.5)--(2.4.7), pp.~17--18;
\sscite{whittle1954stationary}, \S9, pp.~447--448; \sscite{lindgren2011explicit}, \S2.1,
eqs.~(1)--(2), p.~426). Carried to every elliptical aperture by the linear maps of the image, a
rotation-invariant self-decomposable law gives the family called its \emph{affine orbit}. The orbits
of the isotropic stable laws, with transforms $e^{-|\omega|^\alpha}$ of index $0 < \alpha < 2$, cascade
along the changes of shape that are slow enough for the growth in size, and those of the \Matern{}
laws along none (Propositions~\ref{prop:stable-step-cone} and~\ref{prop:matern-steps}).

\paragraph{The classification.} The main result of the paper is the classification of the families
in dimension $d \ge 2$ over the two choices, Table~\ref{tab:landscape} in \S\ref{sec:landscape-table}. It takes three
groups: $\GL(d)$; the similarities $\mathrm S(d)$, the rotations and reflections composed with the
dilations; and the positive diagonal group $\mathrm D_+(d)$, which rescales the axes of a sampling
grid independently. It sets them against three kinds of order: the ray order, the Loewner order and,
between them, the \emph{stable orders} $\le_\alpha$, $0 < \alpha < 2$. Between positive definite
matrices the stages of $\le_\alpha$ are exactly the steps along which the isotropic stable law of index $\alpha$, carried to every shape by
linear maps, cascades (Proposition~\ref{prop:stable-step-cone}). For each group and each order the
table says what is known of the families, and cites the statements that prove it.
% Round 3 (content review), A.R1 and C-R2 with E i (VERIFY-round3-front-classical-paper.md A.R1;
% PLAN-review-response-paper.md R3.1, "taken unless vetoed"): the complete cells named, three since D10.
For the general linear group the description is complete along the rays, along the stable orders and
along the Loewner order. In the other cells it is a bound, examples or a partial description
(\S\ref{sec:conclusion}). The families are monotone
in both choices. A family of a larger group, or along a larger order, restricts to a family of every
smaller pair, with the same apertures (Lemma~\ref{lem:gw-kernels-from-origin}(1)). So the families
grow as the group shrinks and as the order shrinks towards the rays, and several of these inclusions
are strict. The affine Gaussian scale space is the answer for the largest group with the largest
order. The rest of the table says what each relaxation buys, a smaller group or a smaller order. For
shape adaptation its main reading is the price of a kernel other than the Gaussian, kernel by kernel:
along which changes of shape the record stays a further smoothing, and how fast the shape may change
along them. The paper reads the table in three parts.

\paragraph{The general linear group with the Loewner order.} For $d \ge 2$ this pair admits the affine
Gaussian scale spaces alone, on kernels and on operators (Theorems~\ref{thm:affine-rigidity}
and~\ref{thm:affine-rigidity-operators}). On the line the same axioms admit every symmetric
self-decomposable law. One step decides it. The axioms make every aperture a linear image of the
aperture at the identity, which is rotation invariant (Lemma~\ref{lem:affine-reduction}). Call a step
$\Sigma \to \Sigma'$ from a positive definite $\Sigma$ \emph{rank-deficient} when it widens the
aperture in some directions and leaves it unchanged in the others, that is, when $\Sigma' - \Sigma$ is
positive semidefinite, singular and not zero. The increment of one such step is a probability measure
only for the Gaussian (Theorem~\ref{thm:rank-deficient-step}). A linear map carries such a step to one
from the identity, so for the general linear group every order of the scheme that contains one
rank-deficient step admits only Gaussian families
(Proposition~\ref{prop:shape-adapted-apertures}(3)). The theorem takes a positive definite base, where a rank-deficient step is one with $\Sigma' - \Sigma$ singular.
The base must not be $0$: from $\Sigma = 0$ every step $0 \to P$ has the increment $\mu_P$.

Shape adaptation takes such a step. In the iteration of Lindeberg and G\aa{}rding the first stage is
isotropic, at $wI$, and the adapted aperture is $\Sigma_w = wM$, with $M$ the shape matrix that the
second-moment matrix of the image gives at the current pass, divided by its least eigenvalue (\sscite{lindeberg1994shape}, p.~396, step~2;
\sscite{lindeberg1997shape}, \S6.2). So $\Sigma_w$ keeps the least eigenvalue $w$, and for $M \ne I$
the step $\Sigma_w - wI = w(M - I)$ is rank-deficient (Proposition~\ref{prop:shape-adapted-apertures}(1)).
Among the affine orbits of the rotation-invariant laws with a positive transform only the Gaussian
reaches $\Sigma_w$ by a further smoothing of the first stage
(Proposition~\ref{prop:shape-adapted-apertures}(2)). For every other kernel the record at $\Sigma_w$
is a new convolution of the original image, which each pass of the iteration computes in any case
(\S\ref{sec:landscape-uses}). A stable orbit reaches the adapted shape by a further smoothing only with
a wider minor axis, and a \Matern{} orbit not at all (\S\ref{sec:shape-adapted}). So, as far as the apertures of the practice are concerned, a non-Gaussian kernel loses only the
cascade between the stages of the iteration, which the iteration does not use. How the selection of
the shape behaves under such a kernel is not treated (\S\ref{sec:conclusion}).

\paragraph{The general linear group from the rays to the Loewner order.} This part carries the reading
for shape adaptation: the price of a kernel other than the Gaussian, kernel by kernel. Between the rays
and the Loewner order the kernels that keep a cascade decrease. Along the rays alone the families of
the general linear group are the affine orbits of the rotation-invariant self-decomposable laws
% Round 3 (content review), A.S7: Proposition 3.4(3) cited beside 4.5, as Table 2 does.
(Propositions~\ref{prop:gw-ray-families}(3) and~\ref{prop:ray-families}). Read on the multiples of the identity they are, up to a change
of the scale coordinate, the isotropic families of the isotropic paper
(Proposition~\ref{prop:isotropic-corner}). A chain of orders indexed by $0 \le \alpha \le \infty$ runs
from the rays to the Loewner order, through the stable orders $\le_\alpha$, $0 < \alpha < 2$
(Proposition~\ref{prop:stable-order-generated}). Along it the families of the general linear group
shrink from all the affine orbits of the rays to the Gaussian (Theorem~\ref{thm:wedge-axis}). For
$0 < \alpha, \gamma < 2$ the orbit of the isotropic stable law of index $\gamma$ cascades along the
order of index $\alpha$ exactly when $\gamma \ge \alpha$, and from $\alpha = 2$ on only the Gaussian does.
Each order limits how fast the shape may change as the scale grows, and Table~\ref{tab:speed-limit} gives
% Round 3 (content review), A.S14: Table 3 has "at least" and "not known in general" rows.
the limit for each family it lists. The orbits whose transform decays like a power of the frequency keep at
most the rays, so every change of shape is a re-measurement for them
(Proposition~\ref{prop:log-orbit-rays}). The \Matern{} orbits are among them
(Proposition~\ref{prop:matern-steps}). So are the orbits of a law whose transform is a decaying
rational function of $\langle\omega,\Sigma\omega\rangle$ without zeros on $[0,\infty)$
(Proposition~\ref{prop:rational-symbols}). These are the kernels computed by finitely many elliptic
solves in the shape's operator $-\nabla\cdot\Sigma\nabla$.

The orders between the rays and the Loewner order are a question of matrix analysis and of the theory
of Lie semigroups, and from \S\ref{sec:invariant-wedges} on the paper uses their words. An order
invariant under $\GL(d)$ is read off the set of the relative steps $\Sigma^{-1/2}\Sigma'\Sigma^{-1/2}$
that it allows (Proposition~\ref{prop:invariant-orders}). A \emph{cone field} assigns to each $\Sigma$
the cone $\Sigma^{1/2}C\Sigma^{1/2}$ of allowed directions, for a closed convex cone $C$ of symmetric
matrices invariant under rotations and containing no line. It generates the order in which two
matrices are related when a curve that moves inside the cones joins them, as in the theory of Lawson
and of Hilgert, Hofmann and Lawson (\sscite{lawson1989ordered}, Def.~1.6, p.~277, and Thm.~4.9,
p.~290; \sscite{hilgert1989lie}, Def.~VI.1.7, p.~493, and Thm.~VI.1.16, p.~500). These are exactly the closed invariant partial
orders whose steps are closed under fractional powers (Theorem~\ref{thm:cone-orders}), and the order of index
$\alpha$ in the chain above is that of a cone $C_\alpha$ (Proposition~\ref{prop:stable-order-generated}).

\paragraph{Smaller groups.} A smaller group admits more families, and several of the inclusions are
strict (\S\ref{sec:smaller-groups}). Along the rays the isotropic members of every row are the
self-decomposable laws its group allows. Along the general linear row they shrink to the Gaussian,
while the similarity row keeps every rotation-invariant one in every column, the Loewner column included
% Round 3 (content review), A.R9: the witness of F.3(4), the shape-blind family F.3(3)(ii). A.x: the
% forward reference for "exponent". E item 6: the "So" sentence cites H.1(3) as well.
(Proposition~\ref{prop:similarity-ray-families}(4)), by a family that sees only the volume of the
aperture. For the similarities with the Loewner order non-Gaussian families
survive (Proposition~\ref{prop:similarity-stable-family}(3)). The exponent (\S\ref{sec:line-axioms}) of their non-Gaussian part
changes under a rotation by at most a factor of four times the eccentricity of the aperture, the ratio
of its extreme eigenvalues (Proposition~\ref{prop:similarity-stable-order-kernels}(2)). So the rigidity of the general linear
group with the Loewner order needs both: the Loewner order alone does not force the Gaussian, and the
general linear group alone does not either (Proposition~\ref{prop:similarity-stable-family}(3) and
Theorem~\ref{thm:wedge-axis}(3)). For the positive diagonal group a family that cascades along each axis is, at the diagonal
apertures, a product of self-decomposable laws on the line, one for each axis
(Proposition~\ref{prop:fixed-frame-families}). Such families gain anisotropy along the axes as fast as the
Gaussian does and cannot turn. Along the rays a family of any group is a continuous covariant field of
self-decomposable laws over the shapes (Proposition~\ref{prop:gw-ray-families}). The price of a smaller
% Round 3 (content review), E item 7 (VERIFY-round3-smaller-groups-paper.md section 7).
group is shape selection that transforms with the image, which needs a group that carries every shape
to every other (\S\ref{sec:landscape-uses}).

\paragraph{Scope.} The axioms are stated on operators and on their kernels, and for every pair
$(G, W)$ the two forms describe the same families (Lemmas~\ref{lem:operators-to-kernels}
and~\ref{lem:kernels-to-operators}). So each statement on the kernel families of a pair is one on its
operator families. The isotropic theory in $\RR^d$ is that of the isotropic paper, used through the
% Round 3 (content review), A.S5: Props. 3.11, 3.12(2), Thm. 3.13 and Rem. 3.26 of the isotropic paper.
six statements restated in \S\ref{sec:iso-recap}, and five more cited where they are used. The paper is theory and makes no empirical claim
about images. How a shape is selected, how the iteration behaves and singularity theory over the
larger group are not treated, and \S\ref{sec:conclusion} names the limits of the results. The
computation of these families on grids is left to a further paper, in preparation.
% TODO(module VI): cite spatial-hemigroup-pyramids (Paper VI) when released.

\paragraph{Relation to scale-space theory.} The affine Gaussian scale space is classical, and so is
its derivation from axioms. Iijima derived it in 1962 and 1963, as Weickert, Ishikawa and Imiya
report (\sscite{weickert1999linear}, p.~241). His scale is a positive definite $2 \times 2$ matrix,
his transformations are linear, translation invariant, positive and closed under every invertible
linear map, and any two of them compose to a third. Lindeberg obtains the affine Gaussian scale space from a semigroup in a scalar scale with the
non-enhancement of local extrema (\sscite{lindeberg2011generalized}, Thm.~5, p.~42).
Theorem~\ref{thm:affine-rigidity} rests on other axioms. Its cascade is asked for only from a matrix to
a larger one in the Loewner order, and its increments need not belong to the family. It holds for every
$d \ge 2$, and by Theorem~\ref{thm:rank-deficient-step} one rank-deficient step is enough. Neither
earlier result implies it, and it implies neither of them (Remark~\ref{rem:lindeberg}). The morphological axiomatics of Alvarez, Guichard, Lions and Morel keeps
the scale a scalar, and a linear map of the image is absorbed there by a change of that scalar
(\sscite{alvarez1993axioms}, p.~206). The $\alpha$-scale spaces of Pauwels, Van Gool, Fiddelaers and
Moons and of Duits, Florack, de Graaf and ter Haar Romeny admit the stable laws
(\sscite{pauwels1995extended}, p.~694; \sscite{duits2004axioms}, p.~277). The general linear group with
the Loewner order removes them. For $\gamma < 2$ the orders of the chain up to index $\gamma$ keep the
stable law of index $\gamma$ (Theorem~\ref{thm:wedge-axis}). \S\ref{sec:classical-prior-art} sets the
axiomatics side by side.

Anisotropy is often counted among the minor modifications of the Gaussian scale space, those that can
be ``easily accounted for'' (Duits, Felsberg, Florack and Platel, \sscite{duits2003scale}, p.~494).
For a kernel other than the Gaussian, \S\ref{sec:landscape} gives the cost of a change of shape.
In the scale-space and vision work read for this paper, neither the order of a linear family nor its
group is posed as a choice. Lindeberg states the order as the meaning of a larger matrix scale (\S\ref{sec:affine-covariance}),
and in feature detection the group follows the application. Tuytelaars and Mikolajczyk choose between
scale-invariant and affine-invariant detectors by the deformations expected, on experimental evidence
(\sscite{tuytelaars2008local}, \S7.2, p.~258). For nonlinear evolutions the group is chosen. Olver,
Sapiro and Tannenbaum classify the evolutions invariant under a prescribed group
(\sscite{olver1997invariant}, \S4, p.~183), and Sapiro and Tannenbaum note that other groups give other
scale spaces (\sscite{sapiro1993affine}, p.~36).

% Q-0314: pedersen2003cone read at the page (records/affine/READING-pedersen2003cone.md); the
% TODO(source) is closed. Q-0317 (the blind register review, section 1 flag 5): the sentence on
% Barndorff-Nielsen--Pedersen--Sato and the members at a basis of the cone is dropped; the contrast
% stays.
\paragraph{Relation to probability and to matrix analysis.} In probability the core of
Theorem~\ref{thm:affine-rigidity} is Kelker's lemma, that a rotation-invariant law with independent
coordinates is normal (\sscite{kelker1970distribution}, Lemma~5, p.~423). Through the
singular matrix $I - e_1e_1^{\tp}$ the cascade makes the first coordinate of $\mu_I$ independent of the
others. Rotation invariance makes every coordinate so, and Kelker's lemma then makes $\mu_I$ normal
(Remark~\ref{rem:kelker}). Theorem~\ref{thm:rank-deficient-step} reads one rank-deficient step from a
positive definite matrix alone. Its increment is any probability measure, and the transforms at its
two ends need not be those of laws. The affine Gaussian
scale space is also a convolution semigroup over the cone $\PSD(d)$,
$\mu_\Sigma * \mu_{\Sigma'} = \mu_{\Sigma+\Sigma'}$. With the covariance $\Sigma$ at $\Sigma$ it is the
canonical $\PSD(d)$-parameter convolution semigroup of Pedersen and Sato (\sscite{pedersen2003cone},
Def.~2.1 and the paragraph after it, p.~505; \sscite{pedersen2004relations}, Def.~3.5 and Ex.~3.6,
p.~546). Their setting assumes additivity and has no covariance
group. Within it, for $d \ge 2$, finite second moments and a covariance at most $\Sigma$ at each $\Sigma$ force
Gaussian laws (\sscite{pedersen2004relations}, Prop.~5.4, p.~556), and a covariance equal to $\Sigma$
leaves only the canonical semigroup, up to drift (Rem.~5.5, p.~557). That is a Gaussian rigidity under
other hypotheses than those of Theorem~\ref{thm:affine-rigidity}: additivity over the whole cone,
moments and a covariance bound, in place of covariance under $\GL(d)$ and the cascade. The scheme assumes neither additivity nor a semigroup, and for the general linear group with the
Loewner order Theorem~\ref{thm:affine-rigidity} derives the additivity, and the Gaussian with it, from
covariance and the cascade. \S\ref{sec:affine-stable-cone} also uses the operator-self-decomposable
laws of Urbanik, in which the contraction $x \mapsto cx$ of self-decomposability is replaced by the
flow $e^{-tB}$ of a matrix $B$ \sscite{urbanik1972levy}. Its cones of such matrices belong to the
setting of Jurek and Mason (\sscite{jurek1993operator}, p.~85 and Prop.~3.3.6, p.~86). The paragraphs
headed ``Prior art'' in \S\ref{sec:invariant-wedges} and \S\ref{sec:affine-stable-cone} say what is
taken from Lawson, from Hilgert, Hofmann and Lawson and from Jurek and Mason. The \Matern{} orbit is
the anisotropic \Matern{} covariance of spatial statistics, in which the shape matrix is a parameter
of the covariance (\sscite{lindgren2011explicit}, \S3.4, p.~437). No change of that matrix other than a
rescaling is a further smoothing (Proposition~\ref{prop:matern-steps}).

\paragraph{Plan of the paper.} \S\ref{sec:line} restates what is used from the line paper, the cone
paper and the isotropic paper, and the five theorems of the literature that the proofs cite, the
\emph{cited facts}. \S\ref{sec:classical} states the $(G, W)$ scheme and treats the general linear
group with the Loewner order.
On the way it shows that the increments of a continuous cascade along every order with a path condition are
\emph{infinitely divisible}, each
the $n$-fold convolution power of a law for every $n$ (Theorem~\ref{thm:cascade-infinitely-divisible}).
\S\ref{sec:landscape} gives the classification. It says what an affine scale space is for
(\S\ref{sec:landscape-uses}) and states the table of the families of each group along each order
(\S\ref{sec:landscape-table}). Then it treats the invariant orders, the steps an orbit survives, the
general linear group along the chain of cones, shape-adapted smoothing and the smaller groups.
\S\ref{sec:conclusion} is the conclusion: it reads the results for shape adaptation and names their
limits, and \S\ref{sec:open-problems} collects the open problems.
Appendices~\ref{app:classical}--\ref{app:wedge-axis} hold the proofs deferred from
\S\S\ref{sec:classical} and~\ref{sec:landscape}, and the complete proof of
Theorem~\ref{thm:wedge-axis}. Appendices~\ref{app:orientation}--\ref{app:similarity} hold further
results, among them the orientation-preserving group (Proposition~\ref{prop:glplus-families}), the
stable orbits under a tempering (Proposition~\ref{prop:tempered-step-cone}), the negative mass of a
\Matern{} step across shapes (Proposition~\ref{prop:matern-negative-mass}) and the families tied to a
frame of axes (Proposition~\ref{prop:frame-families}). What is machine-checked, and what it assumes, is
said next, and Appendix~\ref{app:verification} holds the detail.

\subsection{What is machine-checked}\label{sec:trust-base}
% trust base: begin
% Pass 6, the author's review of section 1 (2026-10-26; records/affine/REVIEW-author-paper-affine.md,
% comment 2: "most of 1.1 could be moved to an appendix as we did for the isotropic paper. It is
% mainly the first paragraph that is interesting in an introduction"). The subsection is the
% STATEMENT of the trust base (article-kit WRITING.md section 4): what is machine-checked, the four
% cited facts with what each carries beyond its source, the statements proved in the earlier
% papers, and what is computed. Its inventory (what is proved, theorem by theorem; the statements
% that rest on each cited fact and those that rest on none; the standard theorems; the review of the
% formal statements; where the development is, with the Lean names and the stamped #print axioms
% block) is Appendix app:verification, appendix-verification.tex, under a second pair of "trust
% base" markers, moved there verbatim by records/affine/edits_p6_verification_appendix.py. Nothing
% honesty-critical left this subsection. The text kept here is the former opening paragraph, the
% former "What it rests on" up to the sentence on C2, "What is taken from the earlier papers" and
% "What lies outside the formal development", without their headings, and one sentence on the
% review with the pointer. The cited facts are named here as \ledger{} items and \ref'd through
% their theorems, so the ledger comparison of `linkage check` (LINKAGE.md rule 13) reads A2, A3, A5,
% A6 and A9 on both sides.
% The former header comments of this region (Q-0377; CHANGED Q-0281 wave 2, group F; Q-0381) are
% kept with the text they explain: Q-0377's in the appendix, the other two here.
% Round 3 (content review; the author's decisions D10 and D12, PLAN-review-response-paper.md R3.0):
% prop:affine-rigidity-boundary and prop:wedge-axis-profiles were added, and C5 (ledger A2) for
% prop:wedge-axis-profiles(2). Their Lean landed afterwards (D15), so the sentences saying that they were
% proved in print only are gone.
Every definition, lemma, proposition and theorem of \S\S\ref{sec:classical}--\ref{sec:landscape}
and of the appendices has been formalized in the Lean~4 proof assistant, on its mathematical library
Mathlib, and its proof is machine-checked; so have the six statements restated from the isotropic
paper in \S\ref{sec:iso-recap}. This subsection says what the verified results rest on, and it can be
read without the formal development. A proof marked as a sketch shows the decisive steps, and its
complete proof is in the LaTeX sources deposited with this version (Appendix~\ref{app:verification}
lists these statements).

% CHANGED Q-0281 wave 2, group F: C4 added (ledger A9, the author's decision of 2026-10-05).
Five theorems of the literature are taken as given rather than
reproved, four of probability theory and one of matrix analysis. They are the only assumptions of
the results beyond the logical foundations of the proof assistant, and
\S\ref{sec:cited-facts} states each with the page where it was read.
\begin{itemize}
\item[\ledger{A6}] Khintchine's theorem on the limits of null arrays (Theorem~\ref{thm:cited-A6}).
\item[\ledger{A3}] The \Levy--Khintchine representation of the infinitely divisible laws on $\RR^d$:
(i) its existence, (ii) the uniqueness of the triplet and (iii) its converse
(Theorem~\ref{thm:cited-A3}).
\item[\ledger{A5}] The infinite divisibility of operator-self-decomposable laws with a contracting
exponent (Theorem~\ref{thm:cited-A5}).
\item[\ledger{A9}] An inequality between the eigenvalues of a product of two positive definite
matrices and those of its factors (Theorem~\ref{thm:cited-A9}).
\item[\ledger{A2}] The polar criterion, which describes the self-decomposable laws among the
infinitely divisible ones by their \Levy{} measures (Theorem~\ref{thm:cited-A2}).
\end{itemize}
% Q-0381 (the fidelity review's paper-facing point 2, R7 and R18): what C3 carries beyond its pages,
% said item by item as the review found it at the page images (blueprint/AXIOMS.md A5 and
% AXIOMS-verbatim.md hold the transcriptions; Prop. 1.8.1(b) re-read at printed p. 29).
Two of them go beyond the letter of their sources, and \S\ref{sec:cited-facts} says how. C1 is
stated for laws where its source states it for random variables. C3 joins the last sentence of a
theorem of Jurek and Mason to the paragraph of their book that drops its hypothesis on the support of
the law, and it reads two further things beyond those two passages. C4 states less than its source,
the inequalities without the equality of the totals, which is $\det(MN) = \det M\det N$ and is
proved. C2 and C5 are stated as their sources state them. Appendix~\ref{app:verification} lists the statements that rest on
each.

The statements of the line paper restated in
\S\ref{sec:line-axioms}, and the two restated inside the proof of
Proposition~\ref{prop:isotropic-corner}, are proved in its released version, which is cited by
version and says what of it is machine-checked. The formal development
requires Mathlib and the shared library ScaleSpaceCore, which makes no assumption of its own, and no
formal development of the line paper. The formal proofs take two statements of the line paper from
ScaleSpaceCore, where they are proved: the gauge of \prog{V}{Prop.~6.6}, in the proof of
Proposition~\ref{prop:isotropic-corner}, and the necessity half of
Theorem~\ref{thm:line-classification}. ScaleSpaceCore states the latter with the uniqueness clause of
C2 as a hypothesis, which this development discharges by C2(ii). The fact of \prog{V}{Lem.~6.3} is proved inside the proof of
Proposition~\ref{prop:isotropic-corner}. The six
statements restated from the isotropic paper are proved in its released version, which says what of
it is machine-checked, and are machine-checked in the same development, which the two papers share.

Two remarks report computations
(Remarks~\ref{rem:matern-grid} and~\ref{rem:gaussian-matern-step}), and three figures are drawn by
scripts (Figures~\ref{fig:rank-deficient-step}, \ref{fig:shape-plane} and~\ref{fig:cones-plane}). Each says that its script is
deposited with the paper, and the statements are independent of the computed values.

The formal statements were compared with the printed ones by AI agents, in search of one weaker than
printed, true vacuously or true only by a convention of the proof assistant, and the review found
none. The statement on the use of AI at the end of the paper says how the work was done, and
Appendix~\ref{app:verification} says where the formal development is.
% trust base: end
